% MIREX 2026 Audio Key Detection -- extended abstract (2-4 pages) for the K10 submission.
% Plain two-column article; the ISMIR style is not used because its copyright footer names
% the ISMIR proceedings, which this abstract is not part of.
\documentclass[10pt,twocolumn,a4paper]{article}
\usepackage[T1]{fontenc}
\usepackage[utf8]{inputenc}
\usepackage{times}
\usepackage[margin=19mm,top=20mm,bottom=25mm,columnsep=8mm]{geometry}
\usepackage{amsmath,booktabs,graphicx,url,microtype}
\usepackage[hidelinks]{hyperref}
\usepackage[numbers,sort&compress]{natbib}
\newcommand{\ci}[2]{{\footnotesize[#1, #2]}}
\setlength{\parskip}{0pt}
\pagestyle{empty}

\title{\vspace{-8mm}\textbf{A Small Convolutional Key Classifier Trained on Beatport
Annotations\\[2pt] \large MIREX 2026 Audio Key Detection: system description and dataset-usage disclosure}}
\author{John Hurliman\\
Independent researcher, Oakland, CA, USA\\
\texttt{jhurliman@jhurliman.org} \quad \url{https://github.com/jhurliman/jams}}
\date{}

\begin{document}
\maketitle
\thispagestyle{empty}

\begin{abstract}
We submit K10, a 109k-parameter convolutional network that predicts one of 24 global
keys from a log-magnitude constant-Q transform of the whole track. It is trained
only on the \texttt{beatport\_key} annotations distributed through mirdata (the
GiantSteps-MTG-Keys corpus with revised labels), selected by cross-validation inside
that corpus, and run as a single forward pass on CPU. The submission is
inference-only, MIT-licensed, and packaged as one script with bundled weights. This
abstract describes the system, states precisely how GiantSteps Key was used during the
project's development, and reports the development-time GiantSteps Key scores under
both fifth-credit conventions so that MIREX's results can be read against them.
\end{abstract}

\section{System}
\label{sec:system}

\textbf{Input features.} The audio is decoded, downmixed to mono, and resampled to
22{,}050\,Hz. A constant-Q transform with 24 bins per octave is computed over 8 octaves
from C1 with hop 4{,}096 (about 5.4 frames per second), with a 4-semitone margin on
each side that is cropped away after the transform (208 bins computed, the central 192
kept). Magnitudes are compressed with $\log(1+x)$. The full track is one input; nothing
is cropped or windowed at inference time.

\textbf{Architecture.} Three convolutional blocks, each two $3{\times}3$ convolutions
with batch normalization and ELU followed by $2{\times}2$ max pooling (16, 32, and 64
channels), then dropout, then average pooling over \emph{time only}. Pooling over
frequency is deliberately avoided: it would make the network invariant to
transposition, and key identity lives in absolute frequency position. The pooled
$64{\times}24$ map feeds one linear layer with a 24-way softmax over
$\{12\ \text{tonics}\}\times\{\text{major},\text{minor}\}$. The model has 109k
parameters (0.45\,MB). The design follows the convolutional key-classifier family of
Korzeniowski and Widmer~\citep{korzeniowski2018genre}; it is smaller and has no
ensemble.

\textbf{Inference.} One forward pass over the whole-track CQT, argmax over the 24
classes, no test-time augmentation. The output is written as a single tab-separated
line, tonic (sharps only) and mode, e.g.\ \texttt{C\textvisiblespace major} with a tab
separator. The wrapper sets \texttt{torch} to one thread; a two-minute track takes about
1.7\,s wall on one Apple-Silicon CPU core including interpreter start and model load,
with about 390\,MB peak memory. No GPU and no network access are used at prediction
time.

\textbf{Training.} The corpus is mirdata \texttt{beatport\_key}
\citep{faraldo2017thesis}: 1{,}486 Beatport preview excerpts with revised key
annotations, of which 1{,}363 remain after dropping unparseable or non-major/minor
labels and undecodable audio. Augmentation is a pitch shift of $\pm4$ semitones applied
as a roll of the CQT bins with the label transposed accordingly. Optimization is AdamW
with a cosine schedule on about 60-second random crops. Architecture width and the
epoch budget were chosen by 5-fold cross-validation within the training corpus; the
submitted weights are the model trained on the full corpus at the selected budget. The
whole procedure is committed as \texttt{eval/train\_key\_cnn.py} in the repository, and
the run cost about two hours on one A10 GPU.

\section{Dataset usage and development history}
\label{sec:disclosure}

The task rules ask for an accurate account of dataset use. The account below is taken
from the project's public experiment ledger, which records every evaluation with its
date and commit.

\textbf{Training data.} K10 was trained only on \texttt{beatport\_key}. That corpus
contains all 1{,}157 usable GiantSteps-MTG-Keys training identifiers (it is the same
Beatport collection with revised annotations, not an independent set) and has zero
track-identifier overlap with GiantSteps Key~\citep{knees2015giantsteps}; the check is
committed as \texttt{eval/verify\_key\_disjoint.py} and compares Beatport identifiers
from the two mirdata indexes (1{,}486 versus 600, empty intersection). No Isophonics or
other key data was used. The check does not rule out duplicate recordings under
different identifiers.

\textbf{Prior use of GiantSteps Key in this project.} Before K10 existed, the project
evaluated several key systems on the usable 567 of the 604 GiantSteps Key excerpts
(atonal, multi-key, and unparseable labels dropped; 600 had retrievable audio): a
template matcher (Essentia's \texttt{edma} profile~\citep{faraldo2016key,
bogdanov2013essentia}), the self-supervised S-KEY model~\citep{kong2025skey}, madmom's
CNN key processor~\citep{korzeniowski2018genre,bock2016madmom}, and two small
logistic-regression refinements of the template output that were fit on
GiantSteps-MTG-Keys. One earlier mode classifier had been fit on GiantSteps Key
itself; it was identified as contaminated, retired, and replaced by the
GiantSteps-MTG-fit version before any comparison was reported. The decision to build
K10 came from analysing the banked GiantSteps Key predictions of those systems: on
tracks where madmom was right and our fusion system was wrong, most errors were
unreachable by any reranking of the existing candidates, which argued for a stronger
base model rather than a better combiner. In that sense the benchmark informed the
\emph{decision to train a CNN}, though not its architecture, hyperparameters, or
weights. Two implementation bugs were caught by chance-level cross-validation scores
on the training corpus and fixed before any test-set involvement.

\textbf{Test use of GiantSteps Key by K10.} The experiment was pre-registered in the
ledger before training, with a fixed decision rule. All selection (width, epoch budget,
and the choice between the standalone CNN and two fusion variants that reused the
existing combiner) used training-corpus cross-validation only. The locked model was
then evaluated exactly once on the 567-track subset. A post-hoc per-genre breakdown of
that single evaluation was computed afterwards and is descriptive only. No further
training or tuning followed. The weights in this submission are byte-identical to that
evaluated checkpoint (MD5 \texttt{6141daf25376c16a7bc4326b742e4a3c}).

\section{Development-time results on GiantSteps Key}
\label{sec:results}

Table~\ref{tab:gskey} reports the single evaluation, scored two ways. The project's
scorer awards 0.5 for a same-mode key a perfect fifth away in \emph{either} direction,
as in the 2018 definition~\citep{korzeniowski2018genre}; \texttt{mir\_eval}
0.8.2~\citep{raffel2014mireval} awards it only for the ascending fifth. The task
organizers have stated that both conventions will be reported, so both are given.
Exact, relative (0.3), and parallel (0.2) credits are identical under the two
conventions. Confidence intervals are 95\% percentile bootstraps over tracks
(10{,}000 resamples, seed 0).

\begin{table*}[t]
\centering\small
\setlength{\tabcolsep}{6pt}
\begin{tabular}{lccc}
\toprule
system & weighted, sym.\ fifth & weighted, asc.\ fifth & exact \\
\midrule
K10 (this submission) & 0.8321 \ci{0.804}{0.859} & 0.8145 & 0.7795 \\
madmom CNN (default processor) & 0.8328 \ci{0.806}{0.858} & 0.8134 & 0.7725 \\
\midrule
paired K10 $-$ madmom & $-0.0007$ \ci{-0.019}{+0.018} & $+0.0011$ \ci{-0.020}{+0.022} & \\
wins / losses / ties & 41 / 40 / 486 & 38 / 39 / 490 & \\
\bottomrule
\end{tabular}
\caption{GiantSteps Key, 567 usable excerpts, one evaluation of the locked K10 model.
madmom was run with its published weights and default settings on the same tracks.
The paired difference crosses zero under both fifth conventions; no superiority or
equivalence claim is made.}
\label{tab:gskey}
\end{table*}

These numbers are development-time results on a subset the project had already used,
so they should be read as an expectation, not as a prediction of the MIREX score:
MIREX evaluates the full 604-excerpt set with its own scorer, and a subset mean is not
comparable to a full-set mean. In particular, the project earlier attributed a gap
between madmom's published full-set score (0.746) and its subset score here (0.833) to
subset selection; that explanation was withdrawn because it is arithmetically
impossible under its own assumptions ($604\times0.746/567 = 0.795 < 0.833$), and the
published figure describes a single selected model whereas the default processor is
an ensemble. Descriptively, K10's weighted score by genre on the subset ranges from
0.900 (Dubstep) to 0.739 (Electronica/Downtempo), with Trance (0.763) also below the
mean; sustained, modally ambiguous material remains the main failure mode.

\section{Submission contents and reproducibility}
\label{sec:package}

The package contains \texttt{predict\_key.py} (the calling contract is
\texttt{predict\_key.py \%input \%output}; PEP~723 inline metadata pins
\texttt{torch==2.8.*}, \texttt{librosa>=0.10}, \texttt{numpy>=1.26,<2.3}, Python
3.10--3.12, so \texttt{uv run predict\_key.py in.wav out.txt} builds its own
environment), the weights \texttt{key\_cnn\_v1.pt}, a README with the calling format,
runtime, memory, and environment details, and this abstract. Inference reproduces the
training featurization exactly. The training script, the disjointness check, the
experiment ledger with the pre-registration and the single-shot result, and the
per-track development predictions used for Table~\ref{tab:gskey} are in the public
repository. Code and weights are MIT-licensed; the madmom weights used only for the
comparison above are CC~BY-NC-SA and are not part of the submission.

\section*{Acknowledgments and AI use}
Thanks to the MIREX organizers for clarifying the eligibility question in advance.
Claude (Anthropic) assisted with implementation, experiment execution, and drafting
under the author's direction; OpenAI Codex reviewed the accompanying paper's code and
claims. The author is responsible for the submission.

\bibliographystyle{IEEEtran}
\bibliography{references}

\end{document}
